%% Question bank chapter file - generated by bank_tools.py build (do not edit while the bank is still being built:
%% it is regenerated from the per-test data in Question Bank\data). Every question is one bankq block:
%%   \begin{bankq}[id=..., source={...}, paper=.., qno=.., type=..., subtopic={...}, answer={...}, answer_src=key|solved,
%%                 working={...}, difficulty=1-10, diff_parts={K S I C F}, diff_note={...},
%%                 group=<passage id or empty>, also_in={duplicate copies in other tests}]
%% difficulty = min(10, K+S+I+C+F-1): K concepts 1-3, S steps 1-4, I insight 0-2, C computation 0-2, F format 0-1
%%   ...question LaTeX, options, figures (\bfig / \bfigi -> img/)...
%%   \end{bankq}
%% Passages shared by several questions are bankpassage blocks; their questions carry group=<passage id>.
\documentclass[10pt]{article}
\usepackage[a4paper,margin=1.8cm]{geometry}
\usepackage[T1]{fontenc}
\usepackage{amsmath,amssymb,mathtools}
\usepackage{graphicx}
\usepackage[export]{adjustbox}
\usepackage{xcolor,needspace,keyval}
\usepackage{enumitem}
\usepackage{booktabs,array,tabularx}
\usepackage[hidelinks]{hyperref}
\setlength{\parindent}{0pt}
\setlength{\parskip}{3pt}
\makeatletter
\define@key{bq}{id}{\def\bq@id{#1}}\define@key{bq}{source}{\def\bq@src{#1}}\define@key{bq}{paper}{\def\bq@paper{#1}}
\define@key{bq}{qno}{\def\bq@qno{#1}}\define@key{bq}{type}{\def\bq@type{#1}}\define@key{bq}{subtopic}{\def\bq@sub{#1}}
\define@key{bq}{answer}{\def\bq@ans{#1}}\define@key{bq}{answer_src}{\def\bq@asrc{#1}}\define@key{bq}{working}{\def\bq@work{#1}}
\define@key{bq}{group}{\def\bq@group{#1}}\define@key{bq}{also_in}{\def\bq@also{#1}}
\define@key{bq}{difficulty}{\def\bq@diff{#1}}\define@key{bq}{diff_parts}{\def\bq@dparts{#1}}\define@key{bq}{diff_note}{\def\bq@dnote{#1}}
\newenvironment{bankq}[1][]{\def\bq@work{}\def\bq@also{}\def\bq@group{}\def\bq@diff{?}\def\bq@dparts{}\def\bq@dnote{}\setkeys{bq}{#1}%
  \par\needspace{4\baselineskip}\medskip\noindent{\small\textbf{\bq@id}\hfill\color{gray}[\bq@src, Paper \bq@paper, Q\bq@qno]
  \bq@type\ $\cdot$ \bq@sub\ $\cdot$ \textbf{D\bq@diff}}\par\nopagebreak\relax}
  {\par{\small\color{teal}Answer: \textbf{\bq@ans}\ifx\bq@work\@empty\else\ (solved: \bq@work)\fi
  \ifx\bq@also\@empty\else\par\color{gray}Also in: \bq@also\fi}\par}
\define@key{bp}{id}{\def\bp@id{#1}}\define@key{bp}{source}{\def\bp@src{#1}}\define@key{bp}{questions}{\def\bp@qs{#1}}
\newenvironment{bankpassage}[1][]{\setkeys{bp}{#1}\par\needspace{4\baselineskip}\bigskip\noindent
  {\small\textbf{Passage \bp@id}\hfill\color{gray}[\bp@src; questions \bp@qs]}\par\nopagebreak\relax\itshape}{\par}
\makeatother
\newcommand{\bfig}[2]{\par\smallskip{\centering\includegraphics[width=#1pt,max width=0.62\linewidth,max height=0.32\textheight]{img/#2}\par}\smallskip}
\newcommand{\bfigi}[2]{\includegraphics[width=#1pt,max width=\linewidth,max height=0.32\textheight,valign=c]{img/#2}}
\newcommand{\bmissing}[1]{\textit{[figure missing in the source conversion: \detokenize{#1}]}}
\newenvironment{qoptions}{\begin{itemize}[leftmargin=1.8em,itemsep=1pt,topsep=2pt,label={}]}{\end{itemize}}
\newcommand{\opt}[2]{\item[(#1)] #2}
\providecommand{\fwstart}{}\providecommand{\fwend}{}\providecommand{\qfullsplit}{}
\newcolumntype{L}[1]{>{\raggedright\arraybackslash}p{#1}}
\newcolumntype{C}[1]{>{\centering\arraybackslash}p{#1}}
\newcolumntype{R}[1]{>{\raggedleft\arraybackslash}p{#1}}
\providecommand{\sectionhdr}[1]{}\providecommand{\subsectionhdr}[1]{}


\begin{document}
{\Large\bfseries Surface Chemistry --- Question Bank}\par\smallskip
{\small Chemistry chapter C13. Questions from all converted test papers, grouped by source series. Answers: printed key where it matches the paper, otherwise solved (answer\_src=solved, with working). Exact repeats across papers appear once, with also\_in listing the other copies.}\par\medskip\hrule

\section*{Score-I papers}
\begin{bankq}[id=SI02-P1-CHEM-S2-Q06, source={Score-I - Test 02 (31 Dec 2021)}, paper=1, qno=6, type=multi, subtopic={Colloids}, answer={B,C,D}, answer_src=key, difficulty=4, diff_parts={K1 S2 I1 C0 F1}, diff_note={concepts: properties of colloidal solutions | steps: ultramicroscope; Tyndall; electrophoresis; check}]
Choose \textbf{CORRECT} statement regarding properties of colloidal solutions.
\begin{qoptions}
\opt{A}{Ultramicroscope may be used to determine size of colloidal particles}
\opt{B}{A mixture of milk and water appears blue when observed by reflected light}
\opt{C}{Mixture in option (B), appears red when observed by transmitted light.}
\opt{D}{Brownian movement is responsible for the stability of sols}
\end{qoptions}
\end{bankq}
\begin{bankq}[id=SI02-P1-CHEM-S2-Q12, source={Score-I - Test 02 (31 Dec 2021)}, paper=1, qno=12, type=multi, subtopic={Colloids}, answer={A,B,C}, answer_src=key, difficulty=4, diff_parts={K1 S2 I1 C0 F1}, diff_note={concepts: identify solid sols | steps: solid dispersed in solid: gems, ruby glass; cheese is gel; classify}]
Which among the following may be classified as solid sol ?
\begin{qoptions}
\opt{A}{Pumice stone}
\opt{B}{foam rubber}
\opt{C}{Gemstone}
\opt{D}{Whipped cream}
\end{qoptions}
\end{bankq}
\begin{bankq}[id=SI02-P1-CHEM-S3-Q01, source={Score-I - Test 02 (31 Dec 2021)}, paper=1, qno=1, type=numerical, subtopic={Adsorption}, answer={1}, answer_src=key, difficulty=7, diff_parts={K2 S2 I2 C1 F1}, diff_note={concepts: surface coverage, cylindrical molecules adsorbed on a solid | steps: area per molecule lying or standing; number of molecules; fraction of monolayer; evaluate | I: orientation decides area per molecule}]
A cylindrical molecule `X' has dimensions-radius = 0.5\AA\ and height = 4\AA. `X' molecules are adsorbed on a platinum surface of dimensions [20cm $\times$ 30cm] . If adsorption capacity of Pt- surface is [0.1 $\mu$g of X/cm$^2$] and molar mass of `X' is 60g, then maximum $\mu$-moles of `X' that can be adsorbed on given Pt- surface shall be

[Take : $\pi=\dfrac{22}{7}$, $N_A=6\times10^{23}$] (Note : 1 $\mu$ -- mole = $10^{-6}$ moles)
\end{bankq}
\begin{bankq}[id=SI02-P2-CHEM-S1-Q09, source={Score-I - Test 02 (31 Dec 2021)}, paper=2, qno=9, type=single, subtopic={Colloids}, answer={A}, answer_src=key, difficulty=1, diff_parts={K1 S1 I0 C0 F0}, diff_note={concepts: type of emulsion from soap solubility | steps: soluble in water gives oil in water}]
Neutral soaps which are more soluble in water than in oil, form which type of emulsion ?
\begin{qoptions}
\opt{A}{Oil in water}
\opt{B}{Water in oil}
\opt{C}{Water in water}
\opt{D}{None of the above}
\end{qoptions}
\end{bankq}
\begin{bankq}[id=SI02-P2-CHEM-S1-Q11, source={Score-I - Test 02 (31 Dec 2021)}, paper=2, qno=11, type=single, subtopic={Adsorption}, answer={C}, answer_src=key, difficulty=3, diff_parts={K1 S2 I1 C0 F0}, diff_note={concepts: correct statement about adsorption | steps: physical adsorption exothermic; chemisorption also exothermic; check}]
Which among the following is correct statement -
\begin{qoptions}
\opt{A}{Spontaneous adsorption of gases on solid surface is an endothermic process as entropy decreases during adsorption}
\opt{B}{Formation of micelles takes place when temperature is below Kraft Temperature ($T_k$) and concentration is above critical micelle concentration (CMC)}
\opt{C}{A colloid of $\mathrm{Fe(OH)_3}$ is prepared by adding a little excess (required to completely precipitate $\mathrm{Fe^{3+}}$ ions as $\mathrm{Fe(OH)_3}$) of NaOH in $\mathrm{FeCl_3}$ solution. The particles of this sol will move towards anode during electrophoresis.}
\opt{D}{According to Hardy-Schulze rule for the coagulation of $\mathrm{As_2S_3}$ sol flocculation value of $\mathrm{Fe^{3+}}$ ion will be more than $\mathrm{Ba^{2+}}$ or $\mathrm{Na^+}$.}
\end{qoptions}
\end{bankq}
\begin{bankq}[id=SI02-P2-CHEM-S1-Q14, source={Score-I - Test 02 (31 Dec 2021)}, paper=2, qno=14, type=single, subtopic={Catalysis}, answer={A}, answer_src=key, difficulty=1, diff_parts={K1 S1 I0 C0 F0}, diff_note={concepts: promoters and poisons in catalysis | steps: promoters increase activity; poisons decrease}]
Promoters and Poison are the substances used in chemical reaction which
\begin{qoptions}
\opt{A}{Enhance and decrease the activity of catalyst respectively}
\opt{B}{Decrease and enhance the activity of catalyst respectively}
\opt{C}{Do not have any effect on catalyst}
\opt{D}{Can be used in place of catalyst whenever required}
\end{qoptions}
\end{bankq}
\begin{bankq}[id=SI02-P2-CHEM-S2-Q09, source={Score-I - Test 02 (31 Dec 2021)}, paper=2, qno=9, type=numerical, subtopic={Colloids}, answer={0.30}, answer_src=key, difficulty=5, diff_parts={K1 S2 I1 C1 F1}, diff_note={concepts: mass of protective colloid needed from gold number | steps: gold number is mg protecting 10 mL; scale to 100 mL; evaluate}]
Gold number of Haemoglobin is 0.030. Hence 100 ml of gold sol will require how much haemoglobin (in mg) so that gold is not coagulated by 10 ml of 10\% NaCl solution
\end{bankq}
\begin{bankq}[id=SI08-P1-CHEM-S2-Q06, source={Score-I - Test 08 (09 Feb 2022)}, paper=1, qno=6, type=multi, subtopic={Colloids}, answer={A,C,D}, answer_src=key, difficulty=4, diff_parts={K1 S2 I1 C0 F1}, diff_note={concepts: colloidal systems with liquid dispersion medium | steps: sols, emulsions, gels; liquid medium; check each}]
In which of the following colloidal systems the dispersion medium is liquid -
\begin{qoptions}
\opt{A}{sol}
\opt{B}{gel}
\opt{C}{emulsion}
\opt{D}{froath}
\end{qoptions}
\end{bankq}
\begin{bankq}[id=SI09-P1-CHEM-S2-Q07, source={Score-I - Test 09 (16 Feb 2022)}, paper=1, qno=7, type=multi, subtopic={Colloids}, answer={B,C}, answer_src=key, difficulty=3, diff_parts={K1 S1 I1 C0 F1}, diff_note={concepts: statements related to colloids | steps: peptization, micelle formation, Tyndall effect; check each}]
The correct statement(s) related to colloids is(are)
\begin{qoptions}
\opt{A}{The process of precipitating colloidal sol by an electrolyte is called peptization.}
\opt{B}{Colloidal solution freezes at higher temperature than the true solution at the same concentration.}
\opt{C}{Surfactants form micelle above critical micelle concentration (CMC). CMC depends on temperature}
\opt{D}{Micelles are macromolecular colloids.}
\end{qoptions}
\end{bankq}
\begin{bankq}[id=SI10-P1-CHEM-S2-Q03, source={Score-I - Test 10 (24 Feb 2022)}, paper=1, qno=3, type=numerical, subtopic={Adsorption}, answer={3}, answer_src=key, difficulty=4, diff_parts={K1 S2 I0 C1 F1}, diff_note={concepts: Freundlich isotherm read from a log-log graph | steps: intercept 1 gives k=10; slope 0.25 gives x/m = 10 P^0.25; P^0.25 = 0.3 for the given pressure so adsorbed mass 3 g}]
Adsorption of hydrogen gas over palladium surface obey Freundlich's adsorption isotherm as shown in the graph. The mass (in g) of hydrogen gas adsorbed on the surface of 1.0 g palladium at a pressure of $8.1\times10^{-3}$ atm is
\bfig{103.5}{SI10_p1_chem_s2_q03_1.png}
\end{bankq}
\begin{bankq}[id=SI12-P2-CHEM-S3-Q08, source={Score-I - Test 12 (09 Mar 2022)}, paper=2, qno=8, type=matching, subtopic={Colloids}, answer={P$\to$2, Q$\to$3, R$\to$2}, answer_src=key, difficulty=2, diff_parts={K1 S1 I0 C0 F1}, diff_note={concepts: types of colloids | steps: dispersed phase of paints, fog and smoke}]
\begin{tabular}[t]{|p{0.6cm}|p{2.0cm}|p{0.6cm}|p{2.0cm}|}
\hline
 & \centering List-I & & \centering List-II\tabularnewline
\hline
(P) & Paints & (1) & Gas\\
(Q) & Fog & (2) & Solid\\
(R) & Smoke & (3) & Liquid\\
\hline
\end{tabular}
\end{bankq}
\begin{bankq}[id=SI12-P2-CHEM-S3-Q10, source={Score-I - Test 12 (09 Mar 2022)}, paper=2, qno=10, type=matching, subtopic={Catalysis}, answer={P$\to$3, Q$\to$2, R$\to$1}, answer_src=key, difficulty=2, diff_parts={K1 S1 I0 C0 F1}, diff_note={concepts: catalysts in industrial processes | steps: recall the catalysts of Wacker, Deacon and Ostwald processes}]
Match the catalyst with given process.
\begin{center}\footnotesize
\begin{tabular}{|p{0.6cm}|p{2.0cm}|p{0.6cm}|p{3.0cm}|}
\hline
 & \centering List-I\newline Catalyst & & \centering List-II\newline Process\tabularnewline
\hline
(P) & PdCl$_2$ & (1) & Ostwald process\\
(Q) & CuCl$_2$ & (2) & Deacon process\\
(R) & Pt/Rh & (3) & Wacker process\\
\hline
\end{tabular}
\end{center}
\end{bankq}
\begin{bankq}[id=SI17-P1-CHEM-S1-Q03, source={Score-I - Test 17 (17 Apr 2022)}, paper=1, qno=3, type=single, subtopic={Adsorption}, answer={B}, answer_src=key, difficulty=4, diff_parts={K1 S3 I1 C0 F0}, diff_note={concepts: statements on adsorption and absorption | steps: chalk dipped in ink shows both; silica gel adsorbs water; anhydrous CaCl2 absorbs; magneson test; check each}]
Which of the following statement(s) is(are) correct?

(1) When a chalk stick is dipped in ink, both absorption and adsorption takes place and surface of chalk retains the color of ink while the inside in white.

(2) Water vapours are absorbed by silica gel

(3) Water vapours are adsorbed by anhydrous CaCl$_2$

(4) The precipitate of Mg(OH)$_2$ attains pink colour when precipitated in presence of magneson reagent
\begin{qoptions}
\opt{A}{4 only}
\opt{B}{1 Only}
\opt{C}{2, 3 only}
\opt{D}{1, 2, 3 only}
\end{qoptions}
\end{bankq}
\begin{bankq}[id=SI18-P2-CHEM-S1-Q06, source={Score-I - Test 18 (24 Apr 2022)}, paper=2, qno=6, type=multi, subtopic={Colloids}, answer={A,C}, answer_src=key, difficulty=5, diff_parts={K1 S2 I1 C1 F1}, diff_note={concepts: incorrect statements about surface chemistry | steps: conductivity in emulsions; Tyndall effect conditions; other statements; check each}, also_in={JEE Advanced (2021 Series) - Test 01 (24 Jan 2021) P2 Q6}]
Identify the statement(s) which is/are \textbf{NOT CORRECT} with respect to the surface phenomenon.
\begin{qoptions}
\opt{A}{If on adding electrolyte in an emulsion, the conductivity decreases, then it will be an oil in water type emulsion.}
\opt{B}{Tyndall effect is observed when the refractive indices of the dispersed phase and dispersion medium differ largely.}
\opt{C}{Macromolecular colloids are generally lyophobic in nature.}
\opt{D}{Gases which can react with the adsorbents generally show chemisorption.}
\end{qoptions}
\end{bankq}
\begin{bankq}[id=SI24-P1-CHEM-S1-Q03, source={Score-I - Test 24 (11 Dec 2022)}, paper=1, qno=3, type=single, subtopic={Adsorption}, answer={A}, answer_src=key, difficulty=4, diff_parts={K1 S2 I1 C1 F0}, diff_note={concepts: area occupied by one acetic acid molecule on charcoal | steps: moles adsorbed from fall in molarity times volume; number of molecules with Avogadro number; area of charcoal divided by number of molecules}]
1 g of charcoal adsorbs 100 ml of 0.5 M CH$_3$COOH to form a monolayer, and thereby the molarity of CH$_3$COOH reduces to 0.49. Find the surface area of the charcoal adsorbed by each molecule of acetic acid. Surface area of charcoal $=3.01\times10^{2}$m$^2$/g.
\begin{qoptions}
\opt{A}{$5\times10^{-19}$ m$^2$}
\opt{B}{$9\times10^{-19}$ m$^2$}
\opt{C}{$20\times10^{-19}$ m$^2$}
\opt{D}{$25\times10^{-19}$ m$^2$}
\end{qoptions}
\end{bankq}
\begin{bankq}[id=SI24-P1-CHEM-S2-Q09, source={Score-I - Test 24 (11 Dec 2022)}, paper=1, qno=9, type=multi, subtopic={Catalysis}, answer={A,B,C}, answer_src=key, difficulty=5, diff_parts={K1 S3 I1 C0 F1}, diff_note={concepts: properties of zeolites | steps: hydrated aluminosilicates with shape selective catalysis; general formula; pore sizes; ZSM-5 converts alcohols to gasoline; check each}]
Which of following statement(s) is/are correct for Zeolites :
\begin{qoptions}
\opt{A}{They are hydrated alumino-silicates which can be used for shape selective catalysis}
\opt{B}{They have general formula $\mathrm{M_{x/n}[(AlO_2)_x(SiO_2)_y]}$\\$\mathrm{.mH_2O}$ where, n is the charge on the metal cation $\mathrm{M^{n+}}$}
\opt{C}{They have pore sizes between 260 pm to 740 pm}
\opt{D}{ZSM-5 which is used in petroleum industry as a catalyst, converts hydrocarbons into alcohols is an example of Zeolites.}
\end{qoptions}
\end{bankq}
\begin{bankq}[id=SI31-P2-CHEM-S2-Q04, source={Score-I - Test 31 (29 Dec 2022)}, paper=2, qno=4, type=numerical, subtopic={Colloids}, answer={5}, answer_src=solved, working={Correct statements: (a) soap emulsifies, (c) electrophoresis can coagulate a sol, (d) larger zeta potential gives more stable colloid, (f) lower gold number gives more protection, (g) lower coagulation value gives more coagulating power; (b) Bredig arc is for metal sols and (e) methylene blue is a positive sol, so 5. The printed key has no value for this question.}, difficulty=4, diff_parts={K1 S1 I1 C1 F1}, diff_note={concepts: count of correct statements about colloids and soaps | steps: check each of seven statements about emulsifier Bredig arc coagulation zeta potential charge of sols gold number coagulation value; count}]
How many of the following statements is/are correct?

(a) Soap acts as emulsifier in its cleansing action.

(b) Bredig's Arc method is particularly suitable for preparing sulphur sol.

(c) Coagulation of sol can be caused by electrophoresis.

(d) Greater is the zeta potential greater is stability of the colloid.

(e) Methylene blue sol is a negatively charged sol whereas TiO$_2$ sol is a positively charged sol.

(f) Lower the gold number, more is the protective power.

(g) Lower the coagulation value, higher the coagulation power.
\end{bankq}

\section*{Score-II papers}
\begin{bankq}[id=SII01-P2-CHEM-S1-Q01, source={Score-II - Test 01 (21 Mar 2021)}, paper=2, qno=1, type=multi, subtopic={Colloids}, answer={B,C}, answer_src=key, difficulty=4, diff_parts={K1 S2 I1 C0 F1}, diff_note={concepts: electrophoresis prevention phenomena | steps: electro-osmosis reverses movement; coagulation; check statements}]
Which phenomenon occurs when an electric field is applied to a colloidal solution and electrophoresis is prevented by some suitable means?
\begin{qoptions}
\opt{A}{Reverse osmosis takes place.}
\opt{B}{Electrosmosis takes place}
\opt{C}{Dispersion medium begins to move}
\opt{D}{Dispersion medium becomes stationary}
\end{qoptions}
\end{bankq}
\begin{bankq}[id=SII03-P1-CHEM-S1-Q05, source={Score-II - Test 03 (05 May 2021)}, paper=1, qno=5, type=single, subtopic={Colloids}, answer={A}, answer_src=key, difficulty=2, diff_parts={K1 S1 I0 C1 F0}, diff_note={concepts: positively charged sols | steps: hydrated metal oxides are positive; metal sulphides negative}]
In which of the following given sol or sol made in the process is positively charged.
\begin{qoptions}
\opt{A}{$\text{FeCl}_3+\text{H}_2\text{O}\rightarrow\underset{\text{sol}}{\text{Fe(OH)}_3}+\text{H}^++\text{Cl}^-$}
\opt{B}{$\text{As}_2\text{O}_3+\text{H}_2\text{S}\rightarrow\underset{\text{sol}}{\text{As}_2\text{S}_3}+\text{H}_2\text{O}$}
\opt{C}{$\text{AgNO}_3+\underset{\text{excess}}{\text{KI}}\rightarrow\underset{\text{sol}}{\text{AgI}}+\text{KNO}_3\text{(aq)}$}
\opt{D}{Au sol}
\end{qoptions}
\end{bankq}
\begin{bankq}[id=SII03-P1-CHEM-S1-Q06, source={Score-II - Test 03 (05 May 2021)}, paper=1, qno=6, type=single, subtopic={Colloids}, answer={B}, answer_src=key, difficulty=2, diff_parts={K1 S1 I0 C1 F0}, diff_note={concepts: surfactant with the lowest critical micelle concentration | steps: longer hydrocarbon chain lowers CMC; ionic versus nonionic}]
Which of the following surfactant will form micelles in aqueous solution at lowest molar concentration.
\begin{qoptions}
\opt{A}{$\text{CH}_3-\underset{\text{CH}_3}{\text{CH}}-(\text{CH}_2)_{12}-\text{COO}^-\text{Na}^+$}
\opt{B}{$\text{CH}_3-\underset{\text{CH}_3}{\text{CH}}-(\text{CH}_2)_{15}-\text{OSO}_3^-\text{K}^+$}
\opt{C}{$\text{CH}_3-\underset{\text{CH}_3}{\text{CH}}-(\text{CH}_2)_{10}-\text{CH}_2-\text{NH}_3^+\text{Br}^-$}
\opt{D}{$\text{CH}_3-\underset{\text{CH}_3}{\text{CH}}-(\text{CH}_2)_{11}-\text{COO}^-\text{K}^+$}
\end{qoptions}
\end{bankq}
\begin{bankq}[id=SII04-P2-CHEM-S1-Q04, source={Score-II - Test 04 (07 May 2021)}, paper=2, qno=4, type=multi, subtopic={Colloids}, answer={A,B,D}, answer_src=key, difficulty=5, diff_parts={K1 S3 I1 C0 F1}, diff_note={concepts: colloid statements | steps: sol charge; Brownian motion and Tyndall effect; gel definition; number of colloidal particles; check}, also_in={JEE Advanced (2021 Series) - Test 09 (07 May 2021) P2 Q4}]
Which of the following are correct statements?

\begin{qoptions}
\opt{A}{As$_2$S$_3$ is a negatively charged sol.}
\opt{B}{Brownian movement and Tyndall effect are shown by colloids}
\opt{C}{When liquid is dispersed in liquid, it is called gel.}
\opt{D}{Gold number is a measure of protective power of lyophilic colloid.}
\end{qoptions}
\end{bankq}
\begin{bankq}[id=SII06-P2-CHEM-S1-Q06, source={Score-II - Test 06 (05 Jul 2022)}, paper=2, qno=6, type=multi, subtopic={Adsorption}, answer={B}, answer_src=key, difficulty=3, diff_parts={K1 S2 I0 C0 F1}, diff_note={concepts: adsorption and catalysis | steps: Langmuir theory assumptions; effect of pressure on order; temperature effect on desorption; ZSM-5 shape selective catalyst}]
Which of the following is/are correct?
\begin{qoptions}
\opt{A}{Langmuir adsorption theory states that reactivity of a site also depends on the vacancy status of nearby sites.}
\opt{B}{Order of adsorption can be changed by changing pressure conditions.}
\opt{C}{Desorption of gas on solid is favoured by low temperature.}
\opt{D}{ZSM--5 is shape independent catalyst used to crack hydrocarbons.}
\end{qoptions}
\end{bankq}
\begin{bankq}[id=SII07-P1-CHEM-S1-Q09, source={Score-II - Test 07 (10 Jul 2022)}, paper=1, qno=9, type=single, subtopic={Colloids}, answer={C}, answer_src=key, difficulty=1, diff_parts={K1 S1 I0 C0 F0}, diff_note={concepts: seeding of aluminium hydroxide precipitation | steps: freshly prepared aluminium hydroxide used as seed in the Bayer process}]
For inducing the precipitation of Al(OH)$_3$ from a solution of sodium silicate and sodium aluminate, small amount of freshly prepared sample of which substance is added into the solution ?
\begin{qoptions}
\opt{A}{Na$_2$[Si(OH)$_4$]}
\opt{B}{NaOH}
\opt{C}{Hydrated Al$_2$O$_3$}
\opt{D}{SiC}
\end{qoptions}
\end{bankq}
\begin{bankq}[id=SII07-P1-CHEM-S1-Q18, source={Score-II - Test 07 (10 Jul 2022)}, paper=1, qno=18, type=single, subtopic={Catalysis}, answer={C}, answer_src=key, difficulty=2, diff_parts={K1 S2 I0 C0 F0}, diff_note={concepts: criteria of catalysis | steps: catalyst unchanged chemically; small quantity; does not change equilibrium; find the untrue statement}]
Regarding criteria of catalysis which one of the following statement is not true ?
\begin{qoptions}
\opt{A}{The catalyst is unchanged chemically at the end of the chemical reaction.}
\opt{B}{A small quantity of catalyst is often sufficient to bring about a considerable amount of the reaction.}
\opt{C}{In reversible reaction, the catalyst alters the equilibrium position.}
\opt{D}{The catalyst accelerates the rate of reaction.}
\end{qoptions}
\end{bankq}
\begin{bankq}[id=SII08-P1-CHEM-S1-Q02, source={Score-II - Test 08 (05 Feb 2023)}, paper=1, qno=2, type=multi, subtopic={Adsorption}, answer={B,D}, answer_src=key, difficulty=3, diff_parts={K1 S2 I0 C0 F1}, diff_note={concepts: characteristics of adsorption | steps: chemisorption is monolayer and exothermic; physisorption enthalpy 20-40 kJ/mol; low temperature favours physisorption}]
The correct option(s) related to adsorption processes is(are)
\begin{qoptions}
\opt{A}{Chemisorption results in a multimolecular layer.}
\opt{B}{The enthalpy change during physisorption is in the range of 20 to 40 kJ mol$^{-1}$.}
\opt{C}{Chemisorption is an endothermic process.}
\opt{D}{Lowering the temperature favors physisorption processes.}
\end{qoptions}
\end{bankq}
\begin{bankq}[id=SII09-P2-CHEM-S3-Q04, source={Score-II - Test 09 (12 Feb 2023)}, paper=2, qno=4, type=numerical, subtopic={Colloids}, answer={6}, answer_src=key, difficulty=4, diff_parts={K1 S2 I0 C1 F1}, diff_note={concepts: total surface area after cutting a cube | steps: side of each small cube 1e-4 cm; area of one small cube; number 1e12; total area in cm^2 then m^2}]
A cube of 1 cm side when cut into $10^{12}$ equal cubes, The total surface area (in $m^2$) of all the cubes thus obtained is?
\end{bankq}
\begin{bankq}[id=SII12-P2-CHEM-S1-Q02, source={Score-II - Test 12 (04 Feb 2024)}, paper=2, qno=2, type=single, subtopic={Colloids}, answer={C}, answer_src=key, difficulty=2, diff_parts={K1 S2 I0 C0 F0}, diff_note={concepts: properties of colloidal systems | steps: lyophobic colloids are formed by special methods; emulsions are liquid-liquid; Tyndall effect and others; check statements}]
Consider the following statements related to colloids.\\
(I) Lyophobic colloids are \textbf{not} formed by simple mixing of dispersed phase and dispersion medium.\\
(II) For emulsions, both the dispersed phase and the dispersion medium are liquids.\\
(III) Micelles are produced by dissolving a surfactant in any solvent at any temperature.\\
(IV) Tyndall effect can be observed from a colloidal solution with dispersed phase having almost same refractive index as that of the dispersion medium.\\
The option with the incorrect set of statements is
\begin{qoptions}
\opt{A}{(I) and (II)}
\opt{B}{(II) and (III)}
\opt{C}{(III) and (IV)}
\opt{D}{(II) and (IV)}
\end{qoptions}
\end{bankq}
\begin{bankq}[id=SII15-P2-CHEM-S2-Q06, source={Score-II - Test 15 (09 Feb 2024)}, paper=2, qno=6, type=multi, subtopic={Colloids}, answer={D}, answer_src=key, difficulty=4, diff_parts={K1 S2 I1 C0 F1}, diff_note={concepts: colloid of ferric oxide in excess NaOH | steps: negative sol since OH- adsorbed; coagulation by cations; statements}]
A colloid of ferric oxide is prepared by mixing FeCl$_3$ in excess of NaOH solution. Select the correct statement(s) for this colloid :
\begin{qoptions}
\opt{A}{In electro-osmosis, the sol is coagulated at cathode.}
\opt{B}{It can be coagulated by anions.}
\opt{C}{It is mutual coagulated by addition of colloid of As$_2$S$_3$.}
\opt{D}{It is negatively charged sol.}
\end{qoptions}
\end{bankq}
\begin{bankq}[id=SII15-P2-CHEM-S3-Q02, source={Score-II - Test 15 (09 Feb 2024)}, paper=2, qno=2, type=numerical, subtopic={Adsorption}, answer={8}, answer_src=key, difficulty=5, diff_parts={K1 S2 I1 C1 F1}, diff_note={concepts: mass of gas adsorbed from isotherm data | steps: Freundlich x/m = k P^(1/n); determine constants; adsorption at 2 atm over 2 g charcoal}]
In an adsorption experiment,
\par\nobreak
\bfig{137.0}{SII15_p2_chem_s3_q02_1.png}
Calculate the mass of gas (in gram) adsorbed over 2 gram charcoal under a pressure of 2 atm.
\end{bankq}
\begin{bankq}[id=SII21-P1-CHEM-S1-Q02, source={Score-II - Test 21 (14 Mar 2024)}, paper=1, qno=2, type=multi, subtopic={Adsorption}, answer={A,B,D}, answer_src=key, difficulty=3, diff_parts={K1 S2 I0 C0 F1}, diff_note={concepts: characteristics of chemisorption and physisorption | steps: chemisorption gives a monolayer; its enthalpy is 80 to 240 kJ/mol; it is exothermic; low temperature favours physisorption}]
The correct option(s) related to adsorption processes is(are)
\begin{qoptions}
\opt{A}{Chemisorption results in a unimolecular layer.}
\opt{B}{The enthalpy change during chemisorption is in the range of 100 to 140 $kJ\,mol^{-1}$.}
\opt{C}{Chemisorption is an endothermic process.}
\opt{D}{Lowering the temperature favors physisorption processes.}
\end{qoptions}
\end{bankq}
\begin{bankpassage}[id=SII25-P2-CHEM-S2-G9-10, source={Score-II - Test 25 (20 Mar 2024)}, questions={9-10}]
Adsorption is the accumulation of molecular species at the surface. During an experiment CH$_3$COOH molecules were adsorbed on the charcoal having surface area $\sqrt3\times10^3$m$^2$. Assume CH$_3$COOH molecule is spherical with radius 1\AA\ \& whole available charcoal surface is occupied by CH$_3$COOH molecules forming monolayer in two dimensional close packing with maximum efficiency ($N_A=6\times10^{23}$)
\end{bankpassage}
\begin{bankq}[id=SII25-P2-CHEM-S2-Q09, source={Score-II - Test 25 (20 Mar 2024)}, paper=2, qno=9, type=single, subtopic={Adsorption}, answer={B}, answer_src=key, difficulty=4, diff_parts={K1 S2 I1 C1 F0}, diff_note={concepts: moles of acetic acid adsorbed on charcoal | steps: initial and final concentration; volume; moles adsorbed}, group=SII25-P2-CHEM-S2-G9-10]
What will be number of moles of CH$_3$COOH adsorbed on charcoal -
\begin{qoptions}
\opt{A}{1/6}
\opt{B}{1/12}
\opt{C}{2/3}
\opt{D}{1/3}
\end{qoptions}
\end{bankq}
\begin{bankq}[id=SII25-P2-CHEM-S2-Q10, source={Score-II - Test 25 (20 Mar 2024)}, paper=2, qno=10, type=single, subtopic={Adsorption}, answer={B}, answer_src=key, difficulty=4, diff_parts={K1 S2 I1 C1 F0}, diff_note={concepts: percent w/v of initial acetic acid solution | steps: final 1/3 M after adsorption; initial concentration from data; convert to w/v}, group=SII25-P2-CHEM-S2-G9-10]
If after adsorption molarity reduced to 1/3M for CH$_3$COOH solution. Find \%( w/v) for 500ml solution used initially -
\begin{qoptions}
\opt{A}{0.5}
\opt{B}{3}
\opt{C}{4}
\opt{D}{6}
\end{qoptions}
\end{bankq}
\begin{bankpassage}[id=SII33-P2-CHEM-S2-G7-8, source={Score-II - Test 33 (01 Apr 2024)}, questions={7-8}]
Colloids are thermodynamically unstable as it involves creation of surface. Their stability can be attributed to various kinetic factors, like charge etc. Solvation is also another factor that stabilizes lyophilic sols.
\end{bankpassage}
\begin{bankq}[id=SII33-P2-CHEM-S2-Q07, source={Score-II - Test 33 (01 Apr 2024)}, paper=2, qno=7, type=multi, subtopic={Colloids}, answer={A,B,C}, answer_src=key, difficulty=5, diff_parts={K1 S2 I1 C1 F1}, diff_note={concepts: sign of the charge on a sol | steps: AgI colloid in excess KI adsorbs iodide to be negative; SnO2 with NaOH is negative; As2S3 negative; Fe(OH)3 with FeCl3 positive}, group=SII33-P2-CHEM-S2-G7-8]
A negatively charged sol is formed when
\begin{qoptions}
\opt{A}{AgNO$_3$ is treated with excess of KI}
\opt{B}{SnO$_2$ is treated with NaOH}
\opt{C}{H$_2$S is passed into As$_2$O$_3$}
\opt{D}{FeCl$_3$ is added to a fresh precipitate of Fe(OH)$_3$}
\end{qoptions}
\end{bankq}
\begin{bankq}[id=SII33-P2-CHEM-S2-Q08, source={Score-II - Test 33 (01 Apr 2024)}, paper=2, qno=8, type=multi, subtopic={Colloids}, answer={A,C,D}, answer_src=key, difficulty=3, diff_parts={K1 S2 I0 C0 F1}, diff_note={concepts: properties of lyophobic sols | steps: viscosity same as the medium; no extensive hydration; electrophoresis shows migration; particles visible in ultramicroscope}, group=SII33-P2-CHEM-S2-G7-8]
Which of the following are related to lyophobic sols
\begin{qoptions}
\opt{A}{Viscosity is same as that of the medium}
\opt{B}{Extensive hydration takes place}
\opt{C}{Migration of particles under the influence of electric field}
\opt{D}{Very difficult to detect particles even under ultra-microscope}
\end{qoptions}
\end{bankq}
\begin{bankq}[id=SII37-P1-CHEM-S1-Q01, source={Score-II - Test 37 (10 Feb 2025)}, paper=1, qno=1, type=multi, subtopic={Colloids}, answer={BCD}, answer_src=key, difficulty=5, diff_parts={K1 S3 I1 C0 F1}, diff_note={concepts: properties of colloids and surfactants | steps: charge on hydrous ferric oxide sol is positive so it moves to the cathode; gases with higher critical temperature adsorb more; mutual coagulation of oppositely charged colloids; electro-osmosis; CMC trend; Brownian motion and viscosity; test each of the six statements}]
Select correct option among the following statements:

(P) Sol prepared by hydrolysis of FeCl$_3$, colloidal particles get attracted forward anode during electrophoresis.

(Q) A gas with higher critical temperature gets adsorbed to more extent than a gas with lower critical temperature.

(R) On mixing of eosin (acid dye) and haemoglobin (blood), coagulation takes place.

(S) Gold sol is deposited at anode in the electro osmosis process.

(T) CMC value of CH$_3$(CH$_2$)$_8$NH$_3^\oplus$Cl$^\ominus$ will be less than that of CH$_3$(CH$_2$)$_5$COONa.

(U) Brownian motion is independent of nature of colloid but depends on viscosity of solution.

\begin{qoptions}
\opt{A}{PQR}
\opt{B}{QRT}
\opt{C}{RTU}
\opt{D}{QTU}
\end{qoptions}
\end{bankq}
\begin{bankq}[id=SII38-P1-CHEM-S1-Q05, source={Score-II - Test 38 (12 Feb 2025)}, paper=1, qno=5, type=single, subtopic={Colloids}, answer={D}, answer_src=key, difficulty=2, diff_parts={K1 S2 I0 C0 F0}, diff_note={concepts: coagulation of sols by trivalent ions | steps: Al3+ coagulates negatively charged sols; starch is a neutral lyophilic sol; As2S3 and silver sols are negative; methylene blue sol is positive}]
Which of the following sols can not be coagulated by Al$^{3+}$ ion
\begin{qoptions}
\opt{A}{Sols of starch}
\opt{B}{As$_2$S$_3$ sol}
\opt{C}{Silver sol}
\opt{D}{Methylene blue sol}
\end{qoptions}
\end{bankq}
\begin{bankq}[id=SII40-P2-CHEM-S1-Q02, source={Score-II - Test 40 (16 Mar 2025)}, paper=2, qno=2, type=single, subtopic={Colloids}, answer={C}, answer_src=key, difficulty=2, diff_parts={K1 S2 I0 C0 F0}, diff_note={concepts: properties of colloidal systems | steps: lyophobic colloids are formed by special methods; emulsions are liquid-liquid; Tyndall effect and others; check statements}]
Consider the following statements related to colloids.\\
(I) Lyophobic colloids can not be coagulated.\\
(II) For gel, both the dispersed phase and the dispersion medium are liquid.\\
(III) Micelles can not be formed below Kraft temperature.\\
(IV) Tyndall effect can be observed from a colloidal solution with dispersed phase having the different refractive index as that of the dispersion medium.\\
The option with the correct set of statements is
\begin{qoptions}
\opt{A}{(I) and (II)}
\opt{B}{(II) and (III)}
\opt{C}{(III) and (IV)}
\opt{D}{(II) and (IV)}
\end{qoptions}
\end{bankq}
\begin{bankq}[id=SII44-P1-CHEM-S1-Q01, source={Score-II - Test 44 (23 Mar 2025)}, paper=1, qno=1, type=single, subtopic={Colloids}, answer={C}, answer_src=key, difficulty=4, diff_parts={K1 S3 I1 C0 F0}, diff_note={concepts: products of a reaction and properties of a carboxylic acid pair | steps: identify X and Y from the scheme with X heavier; acidity order of the two acids; physical states; use of Y in amino acid synthesis; find the incorrect statement}]
\begin{center}
\bfigi{209.2}{SII44_p2chem_q1_reaction.png}
\end{center}

Given: molar mass of 'X' is greater than 'Y'

'X' and 'Y' are considered as major products

Identify the \textbf{INCORRECT} statement?

\begin{qoptions}
\opt{A}{'X' and 'Y' can be differentiated by NaHCO$_3$ aqeuous solution}
\opt{B}{At room temperature, 'X' is a colourless solid and 'Y' is a colourless liquid}
\opt{C}{Acidic strength: X $>$ Y}
\opt{D}{'Y' can be used to prepare non-essential amino acid}
\end{qoptions}
\end{bankq}

\section*{Score Advanced papers}
\begin{bankq}[id=SA02-P2-CHEM-S1-Q02, source={Score Advanced - Test 02 (13 May 2021)}, paper=2, qno=2, type=multi, subtopic={Colloids}, answer={A,B}, answer_src=solved, working={Gold sol is negative; lyophilic sols act as protective colloids}, difficulty=4, diff_parts={K1 S2 I1 C0 F1}, diff_note={concepts: statements on colloids | steps: charge of sols; protective action; blood}, also_in={JEE Advanced (2021 Series) - Test 10 (13 May 2021) P2 Q2}]
Select the incorrect options :
\begin{qoptions}
\opt{A}{Gold sol is a positive charged sol.}
\opt{B}{Lyophobic sols can act as protective colloid for Lyophilic sols.}
\opt{C}{When FeCl$_3$ is hydrolysed with hot water, then Ferric hydroxide sol (positively charged) will be obtained.}
\opt{D}{Blood is negatively charged sol.}
\end{qoptions}
\end{bankq}
\begin{bankq}[id=SA07-P1-CHEM-S3-Q07, source={Score Advanced - Test 07 (16 Apr 2023)}, paper=1, qno=7, type=numerical, subtopic={Adsorption}, answer={5}, answer_src=key, difficulty=5, diff_parts={K1 S2 I1 C1 F1}, diff_note={concepts: adsorption on charcoal | steps: moles of acid initially; moles after adsorption; Freundlich or mass adsorbed; compute}]
40 ml of 0.5 M aqueous solution of oxalic acid (molar mass = 126) is shaken with 25.2 g of wood charcoal. The final concentration of the solution after adsorption of oxalic acid is 0.25 M. If the weight of oxalic acid adsorbed per g of carbon is represented as y $\times$ 10$^{-2}$ then find the value of y. (Assume volume of oxalic acid remains constant)
\end{bankq}
\begin{bankq}[id=SA09-P2-CHEM-S1-Q06, source={Score Advanced - Test 09 (29 Apr 2023)}, paper=2, qno=6, type=multi, subtopic={Colloids}, answer={B,D}, answer_src=key, difficulty=4, diff_parts={K1 S2 I0 C1 F1}, diff_note={concepts: statements on colloids and gold number | steps: lower gold number means better protective power; electrolytes coagulate oppositely charged sol; impurities destabilise; check statements}, also_in={Pre-Score (2023-24) - Test 07 (29 Nov 2023) P1 Q1}]
Which of the following statement is/are CORRECT?
\begin{qoptions}
\opt{A}{Higher is the gold number of a lyophilic sol. more is its coagulating power}
\opt{B}{An electrolyte can coagulate both positive sol and negative sol}
\opt{C}{Impurities present in a colloid makes it more stable}
\opt{D}{Dialysis is a process to remove impurities of ions and smaller molecules from a colloid}
\end{qoptions}
\end{bankq}
\begin{bankq}[id=SA12-P1-CHEM-S3-Q07, source={Score Advanced - Test 12 (11 May 2023)}, paper=1, qno=7, type=numerical, subtopic={Colloids}, answer={6}, answer_src=key, difficulty=2, diff_parts={K1 S1 I0 C0 F1}, diff_note={concepts: colloids statements | steps: stability of lyophilic; Tyndall}, also_in={Score Advanced - Test 29 (01 May 2024) P2 Q7}]
From the following choose the number of correct statement(s) :-\par
(i) Stability of lyophillic colloids is mainly due to the strong interaction between dispersed phase and dispersion medium.\par
(ii) Entropy change for micelle formation in aqueous solution is positive.\par
(iii) Gelatin has considerably low value of gold number and is an effective protective agent.\par
(iv) Zeta potential is also responsible for stability of lyophobic colloids.\par
(v) Surface tension of lyophilic colloids is lesser than that of the dispersion medium.\par
(vi) A positively charged ion can cause the precipitation of As$_2$S$_3$ sol\par
(vii) Cellulose solution is an example of multimolecular colloid system.\par
(viii) The combination of two layers of same charge around the colloidal particle is called Helmholtz electrical double layer.
\end{bankq}
\begin{bankq}[id=SA13-P1-CHEM-S2-Q05, source={Score Advanced - Test 13 (13 May 2023)}, paper=1, qno=5, type=single, subtopic={Colloids}, answer={5}, answer_src=key, difficulty=3, diff_parts={K1 S3 I0 C0 F0}, diff_note={concepts: colloid and surface chemistry statements | steps: coagulation at delta; second; third; fourth; fifth; count}]
How many of the following statements are correct ?\par
(A) Delta at the mouth of river is formed due to coagulation of the colloidal sea water by the electrolytes present in the river water.\par
(B) Artificial rain is caused by the coagulation of colloidal water droplets in the cloud by an electrolyte.\par
(C) Cottrell precipitator is used in the precipitation of the colloidal waste in the sewage water.\par
(D) Cheese is an example of solid foam.\par
(E) A colloid formed by adding small amount of AgNO$_3$ to KI solution is negatively charged.\par
(F) Micelles are formed only below the Kraft temperature.\par
(G) Mixing of hydrated ferric oxide sol and arsenious sulphide sol bring mutual coagulation.\par
(H) An oil in water emulsion can be diluted by addition of water.\par
(I) Emulsions show Brownian motion and Tyndall effect.\par
\end{bankq}
\begin{bankq}[id=SA16-P2-CHEM-S4-Q09, source={Score Advanced - Test 16 (20 May 2023)}, paper=2, qno=9, type=numerical, subtopic={Colloids}, answer={4}, answer_src=key, difficulty=5, diff_parts={K1 S3 I0 C1 F1}, diff_note={concepts: sols that coagulate haemoglobin | steps: haemoglobin sol is negatively charged; positively charged sols coagulate it: Fe(OH)3 sol; clay and gold and As2S3 are negative; Ca(OH)2 is an electrolyte; count}]
From the given list of sols how many can be used to coagulate the haemoglobin sol by addition to it ?\par
(i) Fe(OH)$_3$ sol\quad (ii) Ca(OH)$_2$ sol\quad (iii) Gold sol\quad (iv) Clay\par
(v) As$_2$S$_3$ sol\quad (vi) Al(OH)$_3$ sol\quad (vii) CdS sol\quad (viii) Cd(OH)$_2$ sol.
\end{bankq}
\begin{bankq}[id=SA17-P1-CHEM-S1-Q03, source={Score Advanced - Test 17 (25 May 2023)}, paper=1, qno=3, type=multi, subtopic={Colloids}, answer={A,C,D}, answer_src=key, difficulty=5, diff_parts={K1 S2 I1 C1 F1}, diff_note={concepts: colloids statements | steps: smoke and cloud both aerosol but different dispersed phase; surface tension of lyophilic sols; Tyndall effect and others; check four}, also_in={Score Advanced - Test 35 (15 May 2024) P1 Q3; Pre-Score (2023-24) - Test 07 (29 Nov 2023) P1 Q5}]
Which of the following statement(s) is(are) correct ?
\begin{qoptions}
\opt{A}{Smoke and cloud both are aerosol with same dispersion medium.}
\opt{B}{Surface tension of lyophilic sols is less than the surface tension of dispersion medium while that of lyophobic sols is more than dispersion medium.}
\opt{C}{Lyophobic sols can be detected by ultra microscope}
\opt{D}{Colloidal sol shows low magnitude of colligative properties, as compared to true solutions.}
\end{qoptions}
\end{bankq}
\begin{bankq}[id=SA19-P1-CHEM-S3-Q06, source={Score Advanced - Test 19 (04 Apr 2024)}, paper=1, qno=6, type=numerical, subtopic={Colloids}, answer={5}, answer_src=key, difficulty=4, diff_parts={K1 S2 I0 C1 F1}, diff_note={concepts: count negatively charged colloids | steps: As2S3 gold clay and acidic dye sols are negative; Fe(OH)3 positive; basic dye positive; count}]
Number of negatively charged colloids among the following is -\par
(i) As$_2$S$_3$ sol\quad (ii) Ferric hydroxide sol\par
(iii) Gold sol\quad (iv) Clay\par
(v) Acidic dye\quad (vi) Hemoglobin\par
(vii) Basic dye\quad (viii) Starch
\end{bankq}
\begin{bankq}[id=SA23-P2-CHEM-S1-Q01, source={Score Advanced - Test 23 (14 Apr 2024)}, paper=2, qno=1, type=multi, subtopic={Colloids}, answer={A,B}, answer_src=solved, working={Emulsions can be diluted with the dispersion medium (A) and broken by heating, freezing or centrifuging (B). Proteins, gums and soaps stabilise O/W emulsions, not W/O (C false). Butter is W/O but cream is O/W (D false).}, difficulty=4, diff_parts={K1 S2 I0 C1 F1}, diff_note={concepts: statements on emulsions and colloids | steps: emulsion can be diluted with the dispersion medium; breaking by heating or freezing; types O/W and W/O; check}]
Select the correct options
\begin{qoptions}
\opt{A}{Emulsion can be diluted with any amount of the dispersion medium}
\opt{B}{Emulsions can be broken into constituent liquid by heating, freezing, centrifuging etc.}
\opt{C}{Emulsifying agents for W/O emulsion are proteins, gums, natural and synthetic soaps etc}
\opt{D}{Butter and Cream are W/O type emulsions}
\end{qoptions}
\end{bankq}
\begin{bankq}[id=SA27-P1-CHEM-S2-Q12, source={Score Advanced - Test 27 (24 Apr 2024)}, paper=1, qno=12, type=multi, subtopic={Adsorption}, answer={A,B}, answer_src=key, difficulty=5, diff_parts={K1 S2 I1 C1 F1}, diff_note={concepts: adsorption graphs | steps: Freundlich isotherm log(x/m) vs log P linear; Langmuir; temperature dependence of adsorption of gas; check each graph}]
Which of the following graphs regarding adsorption are correct?
\begin{qoptions}
\opt{A}{\bfigi{104.8}{SA27_p1_chem_s2_q12_A.png}}
\opt{B}{\bfigi{133.4}{SA27_p1_chem_s2_q12_B.png}}
\opt{C}{\bfigi{130.8}{SA27_p1_chem_s2_q12_C.png}}
\opt{D}{\bfigi{124.9}{SA27_p1_chem_s2_q12_D.png}}
\end{qoptions}
\end{bankq}
\begin{bankq}[id=SA38-P2-CHEM-S4-Q05, source={Score Advanced - Test 38 (13 Apr 2025)}, paper=2, qno=5, type=numerical, subtopic={Adsorption}, answer={1250}, answer_src=key, difficulty=6, diff_parts={K2 S2 I1 C1 F1}, diff_note={concepts: surface area or moles adsorbed from the acetic acid used, monolayer formation | steps: moles taken 100 x 0.6 = 60 mmol; unadsorbed neutralised by 40 mmol NaOH; adsorbed 20 mmol per g; multiply with the area per molecule and Avogadro number}]
To form a complete monolayer of acetic acid on 1g of charcoal, 100 mL of 0.6 M acetic acid was used. Some of the acetic acid ramained unadsorbed. To neutralize the unadsorbed acetic acid, 40 mL of 1 M NaOH solution was required. If each molecule of acetic acid occupies \textbf P $\times$ 10$^{-23}$ m$^2$ surface area on charcoal, the value of \textbf P is \underline{\hspace{1.5cm}}.\par
[Use given data : Surface area of charcoal = 1.5 $\times$ 10$^2$ m$^2$g$^{-1}$; Avogadro's number (N$_\mathrm A$) = 6.0 $\times$ 10$^{23}$ mol$^{-1}$]
\end{bankq}
\begin{bankq}[id=SA39-P1-CHEM-S2-Q10, source={Score Advanced - Test 39 (20 Apr 2025)}, paper=1, qno=10, type=multi, subtopic={Colloids}, answer={A,B,C}, answer_src=key, difficulty=6, diff_parts={K1 S3 I1 C1 F1}, diff_note={concepts: incorrect statements on colloids | steps: collodion is a 4 percent solution of nitrocellulose in ether and alcohol; milk and water look blue by reflected and red by transmitted light; starch gum clay charges; increased ionic strength compresses the double layer; check}]
The incorrect statement(s) related to colloids is/are.
\begin{qoptions}
\opt{A}{The colloidion is a 4\% solution of nitro-cellulose in a mixture of alcohol and ester.}
\opt{B}{A mixture of milk and water appears red when viewed by the reflected light and blue when viewed by the transmitted light.}
\opt{C}{Sols of starch, gum, clay and methylene blue are negatively charged colloidal sol.}
\opt{D}{When the ionic strength of a colloidal solution is increased, thickness of double layer is decreased and the colloid gets precipitated.}
\end{qoptions}
\end{bankq}
\begin{bankq}[id=SA40-P2-CHEM-S4-Q01, source={Score Advanced - Test 40 (27 Apr 2025)}, paper=2, qno=1, type=single, subtopic={Adsorption}, answer={5}, answer_src=key, difficulty=5, diff_parts={K2 S2 I1 C1 F0}, diff_note={concepts: surface area or moles adsorbed from the acetic acid used, monolayer formation | steps: moles taken 100 x 0.6 = 60 mmol; unadsorbed neutralised by 40 mmol NaOH; adsorbed 20 mmol per g; multiply with the area per molecule and Avogadro number}]
To form a complete monolayer of acetic acid on 1g of charcoal, 100 mL of 0.5 M acetic acid was used. Some of the acetic acid remained unabsorbed. To neutralize the unabsorbed acetic acid, 40 mL of 1 M NaOH solution was required. If each molecule of acetic acid occupies \textbf{P} $\times$ 10$^{-21}$ m$^2$ surface area on charcoal, the value of $\left(\dfrac{\mathrm P}5\right)$ is \underline{\hspace{1.2cm}}.\par
[\textbf{Use given data:} Surface area of charcoal = 1.5 $\times$ 10$^2$ m$^2$g$^{-1}$; Avogadro's number (N$_\mathrm A$) = 6.0 $\times$ 10$^{23}$ mol$^{-1}$]
\end{bankq}

\section*{SRG (Enthusiast) papers}
\begin{bankq}[id=SRG16-P1-CHEM-S2-Q07, source={SRG (Enthusiast) - Test 16 (12 Jan 2022)}, paper=1, qno=7, type=multi, subtopic={Colloids}, answer={A,B}, answer_src=solved, working={Lyophilic sols have lower surface tension and higher viscosity than water; colloidal silver and antimony are lyophobic.}, difficulty=2, diff_parts={K1 S1 I0 C0 F1}, diff_note={concepts: properties of lyophilic sols | steps: lyophilic sols have lower surface tension and higher viscosity than water; silver and antimony sols are lyophobic}]
Which of the following is / are correct for lyophilic aquasols ?
\begin{qoptions}
\opt{A}{Its surface tension is lower than that of $\mathrm{H_2O}$}
\opt{B}{Its viscosity is higher than that of water}
\opt{C}{Colloidal silver}
\opt{D}{Colloidal antimony}
\end{qoptions}
\end{bankq}
\begin{bankq}[id=SRG16-P2-CHEM-S3-Q14, source={SRG (Enthusiast) - Test 16 (12 Jan 2022)}, paper=2, qno=14, type=numerical, subtopic={Adsorption}, answer={20}, answer_src=solved, working={n cuts give 8 to the n cubes of edge a/2 to the n, so the area rises by 2 to the n = 1048576 = 2 to the 20.}, difficulty=5, diff_parts={K1 S2 I1 C1 F1}, diff_note={concepts: surface area on repeated cutting of a cube | steps: each cut gives 8 cubes of half the edge; total area increases by 2 each time; solve for n | I: finding the factor of increase of area per cut | C: logarithm base 2}]
A metal cube of edge length `$a$' cm is cut $n$ number of times. It is cut such that, per each cut, each cube is divided into 8 smaller cubes of same edge length. Then what should be the value of $n$, if the final surface area becomes 1048576 times of initial surface area
\end{bankq}
\begin{bankq}[id=SRG16-P2-CHEM-S3-Q17, source={SRG (Enthusiast) - Test 16 (12 Jan 2022)}, paper=2, qno=17, type=numerical, subtopic={Colloids}, answer={50}, answer_src=solved, working={Coagulation value is millimoles of electrolyte per litre of sol: 10 mmol per 0.2 L = 50.}, difficulty=3, diff_parts={K1 S1 I0 C1 F1}, diff_note={concepts: coagulation value of an electrolyte | steps: flocculation value = millimoles per litre of sol | C: unit conversion}]
For the coagulation on 200 mL of $\mathrm{As_2S_3}$ solution, 10 mL of 1 M NaCl is required. What is the coagulating value of NaCl.
\end{bankq}

\section*{Pre-Score (2023-24) papers}
\begin{bankq}[id=PS07-P1-CHEM-S2-Q10, source={Pre-Score (2023-24) - Test 07 (29 Nov 2023)}, paper=1, qno=10, type=single, subtopic={Colloids}, answer={D}, answer_src=key, working={Printed key.}, difficulty=5, diff_parts={K1 S3 I1 C1 F0}, diff_note={concepts: matching colloid terms | steps: precipitation of sol and peptisation; Tyndall; Brownian; dialysis; match | I: recall of definitions | C: four matches}]
Which of the following option show correct matches :-
\par\nobreak
{\small\setlength{\tabcolsep}{3pt}\renewcommand{\arraystretch}{1.3}
\begin{tabular}{|L{0.8cm}|L{3.0cm}|L{0.8cm}|L{2.7cm}|}\hline
\multicolumn{2}{|c|}{\textbf{List-I}} & \multicolumn{2}{c|}{\textbf{List-II}}\\\hline
(P) & Precipitation of colloidal sol & (1) & Peptization\\\hline
(Q) & Purification of colloid & (2) & Coagulation\\\hline
(R) & Gold Number & (3) & Desorption\\\hline
(S) & Formation of colloidal sol & (4) & Dialysis\\\hline
 & & (5) & Measure of protective power\\\hline
\end{tabular}}
\begin{qoptions}
\opt{A}{P $\to$ 1; Q $\to$ 4; R $\to$ 3; S $\to$ 1}
\opt{B}{P $\to$ 5; Q $\to$ 2; R $\to$ 1; S $\to$ 3}
\opt{C}{P $\to$ 2; Q $\to$ 4; R $\to$ 3; S $\to$ 1}
\opt{D}{P $\to$ 2; Q $\to$ 4; R $\to$ 5; S $\to$ 1}
\end{qoptions}
\end{bankq}
\begin{bankq}[id=PS07-P1-CHEM-S3-Q01, source={Pre-Score (2023-24) - Test 07 (29 Nov 2023)}, paper=1, qno=1, type=numerical, subtopic={Colloids}, answer={0.25}, answer_src=key, working={Printed key.}, difficulty=4, diff_parts={K1 S1 I1 C1 F1}, diff_note={concepts: osmotic pressure and statements about sols and solutions | steps: colloidal rubber has far fewer particles than sucrose of same mass so osmotic pressure is smaller; count correct statements | I: colligative property depends on number of particles | C: several statements}]
Consider the following statements:\\
1. Osmotic pressure of rubber sol will be same as that of sucrose solution having same mass of rubber and sucrose mixed in same mass of H$_2$O\\
2. Soap sol of sodium palmitate will be coagulated near cathode on electrophoresis.\\
3. For miscellisation of surfactants entropy of system increases.\\
4. Movement of dispersion medium under the influence of an electric field is called cataphoresis.\\
5. Gold sol is a negatively charged hydrophilic sol.\\
6. Tyndall effect is more distinct for lyophobic colloid than for lyophilic colloids.\\
7. Gold sol on mixing with starch sol causes stablisation of starch sol.\\
8. During adsorption there will be an increase in residual forces on the molecules.\\
9. The extent of adsorption per unit mass of adsorbent will remain same even on changing surface area of same mass of adsorbent.\\
10. On adding 0.1 M KI solution slowly to 0.01 M AgNO$_3$ solution, causes conductivity to increase gradually and after some time it becomes constant.\\
If number of correct statements $=x$\\
And number of incorrect statements $=y$\\
Then, $\dfrac xy=$
\end{bankq}
\begin{bankq}[id=PS07-P1-CHEM-S3-Q08, source={Pre-Score (2023-24) - Test 07 (29 Nov 2023)}, paper=1, qno=8, type=numerical, subtopic={Adsorption}, answer={3.00}, answer_src=key, working={Printed key.}, difficulty=5, diff_parts={K1 S2 I1 C1 F1}, diff_note={concepts: fraction of catalyst sites occupied by hydrogen from desorbed gas | steps: moles of gas from pressure volume and temperature; sites occupied are 10 percent; total sites; compute the number of sites per unit or the surface area | I: moles desorbed correspond to occupied sites | C: gas law conversion and ratio}]
10\% sites of catalyst bed has been occupied by H$_2$ gas molecules. On heating H$_2$ gas is evolved from sites and collected at 0.030 atm and 300K in a vessel of volume 2.46 cm$^3$. Total number of sites available on catalyst is $6\times10^{15}$ per cm$^2$ and surface area of catalyst is $10^3$cm$^2$. The number of H$_2$ molecules adsorbed per site of catalyst are :\\
(Given $N_A=6\times10^{23}$, $R=0.082$ L.atm K$^{-1}$ mol$^{-1}$)
\end{bankq}
\begin{bankq}[id=PS07-P2-CHEM-S1-Q04, source={Pre-Score (2023-24) - Test 07 (29 Nov 2023)}, paper=2, qno=4, type=single, subtopic={Colloids}, answer={C}, answer_src=solved, working={Dispersed phase gas in liquid medium is a foam: whipped cream.}, difficulty=1, diff_parts={K1 S1 I0 C0 F0}, diff_note={concepts: colloids by phase | steps: gas dispersed in liquid is a foam; whipped cream is foam}]
Which of the following is colloid in which dispersed phase is gas and dispersion medium is liquid ?
\begin{qoptions}
\opt{A}{Mist}
\opt{B}{Ruby Glass}
\opt{C}{Whipped Cream}
\opt{D}{Lava}
\end{qoptions}
\end{bankq}
\begin{bankq}[id=PS07-P2-CHEM-S2-Q09, source={Pre-Score (2023-24) - Test 07 (29 Nov 2023)}, paper=2, qno=9, type=multi, subtopic={Catalysis}, answer={B,D}, answer_src=key, working={Printed key.}, difficulty=5, diff_parts={K1 S3 I1 C0 F1}, diff_note={concepts: properties of enzyme catalysed reactions | steps: enzymes are specific; optimum temperature and pH; active site; inhibitors; identify incorrect statements | I: recall of enzyme properties}]
Which of the following statement(s) is/are incorrect with respect to enzyme catalysed reactions?
\begin{qoptions}
\opt{A}{The cavities on enzyme ``Maltase'' \& structure of molecules of Maltose have complimentary shapes for conversion of Maltose to glucose.}
\opt{B}{Just like other catalysts, catalytic action of enzyme is more at a lower temperature.}
\opt{C}{In presence of vitamins \& some metal ions; enzyme activity generally increases.}
\opt{D}{Most of the enzyme catalysed reactions are elementary.}
\end{qoptions}
\end{bankq}
\begin{bankq}[id=PS07-P2-CHEM-S2-Q10, source={Pre-Score (2023-24) - Test 07 (29 Nov 2023)}, paper=2, qno=10, type=multi, subtopic={Colloids}, answer={A,B,C}, answer_src=key, working={Printed key.}, difficulty=5, diff_parts={K1 S2 I1 C1 F1}, diff_note={concepts: coagulation values and the charge of the colloid | steps: coagulating power is inverse of coagulation value; order means the sol is positive when anions are effective; electroosmosis and electrophoresis direction; check | I: coagulation order reveals sign of the charge | C: reading the direction of motion}]
For the coagulation of a particular colloidal solution the order of coagulation value of electrolytes is as follows.
\[\mathrm{Na_3PO_4>BaSO_4>AlCl_3}\]
\begin{qoptions}
\opt{A}{In electroosmosis of this colloidal solution dispersion medium move towards cathode.}
\opt{B}{Diffused layer of the colloidal sol particles contains positively charged ions.}
\opt{C}{It may be a metal sol such as Pt, Au or Hg sol}
\opt{D}{It may be basic dye sol such as methylene blue.}
\end{qoptions}
\end{bankq}
\begin{bankq}[id=PS07-P2-CHEM-S3-Q06, source={Pre-Score (2023-24) - Test 07 (29 Nov 2023)}, paper=2, qno=6, type=numerical, subtopic={Colloids}, answer={6}, answer_src=key, working={Printed key.}, difficulty=5, diff_parts={K1 S2 I1 C1 F1}, diff_note={concepts: sols coagulated by trivalent aluminium | steps: oppositely charged sols are coagulated by the ion; negative sols like Sb2S3 Congo red are coagulated; positive sols like Al2O3 and methylene blue are not; count | I: classify sols by their charge | C: ten sols to classify}]
Total number of sols among below given options which can be coagulated by Al$^{+3}$ ion are?\\
Eosin, methylene blue, Congo red, TiO$_2$, Ag Sol, Sb$_2$S$_3$ Sol, Al$_2$O$_3$ sol. Haemoglobin sol, Clay sol., Charcoal sol.
\end{bankq}
\begin{bankq}[id=PS07-P2-CHEM-S3-Q08, source={Pre-Score (2023-24) - Test 07 (29 Nov 2023)}, paper=2, qno=8, type=numerical, subtopic={Adsorption}, answer={5}, answer_src=key, working={Printed key.}, difficulty=5, diff_parts={K1 S2 I1 C1 F1}, diff_note={concepts: adsorption of acetic acid on charcoal and surface area per molecule | steps: moles adsorbed from the drop in concentration; molecules adsorbed; area per molecule = surface area over number of molecules; scale as needed | I: concentration change gives amount adsorbed | C: unit conversions and large powers}]
1 gm of charcoal having surface area $3.01\times10^2$ m$^2$ is mixed in 100 ml of 0.5 M acetic acid solution; after sufficient time concentration of solution becomes 0.4 M. The surface area occupied by one molecule of acetic acid is $y\times10^{-20}$ m$^2$. Then $y$ is. [Given : $N_A=6.02\times10^{23}$]
\end{bankq}
\begin{bankq}[id=PS08-P1-CHEM-S2-Q06, source={Pre-Score (2023-24) - Test 08 (03 Dec 2023)}, paper=1, qno=6, type=multi, subtopic={Colloids}, answer={C,D}, answer_src=key, working={Printed key.}, difficulty=3, diff_parts={K1 S1 I1 C0 F1}, diff_note={concepts: preparation of colloidal gold | steps: gold sol is prepared by reduction of gold chloride, by Bredig arc method and not by peptisation or hydrolysis | I: classify methods of sol preparation}]
Gold sol can be prepared by :
\begin{qoptions}
\opt{A}{Peptization}
\opt{B}{By hydrolysis}
\opt{C}{Bredic's arc method}
\opt{D}{Reduction}
\end{qoptions}
\end{bankq}
\begin{bankq}[id=PS08-P2-CHEM-S1-Q02, source={Pre-Score (2023-24) - Test 08 (03 Dec 2023)}, paper=2, qno=2, type=single, subtopic={Adsorption}, answer={B}, answer_src=key, working={Printed key.}, difficulty=3, diff_parts={K1 S1 I1 C1 F0}, diff_note={concepts: Freundlich isotherm from a straight line plot | steps: plot of log x/m versus log P has slope 1/n = tan 60 and intercept log k = 0.3010 so k = 2; x/m at 1 atm equals k | I: slope and intercept of the log plot give n and k | C: evaluating with the given numbers}]
A plot for $\log\dfrac{x}{m}$ vs $\log P$ shows a straight line with an angle $60^\circ$ and intercept of 0.3010 on y-axis. The amount of gas adsorbed per gram of adsorbent if initial pressure maintained is 1 atm is:
\begin{qoptions}
\opt{A}{1 gm}
\opt{B}{2 gm}
\opt{C}{3 gm}
\opt{D}{4 gm}
\end{qoptions}
\end{bankq}

\section*{JEE Advanced (2021 Series) papers}
\begin{bankq}[id=JA03-P2-CHEM-S3-Q05, source={JEE Advanced (2021 Series) - Test 03 (01 Mar 2021)}, paper=2, qno=5, type=single, subtopic={Adsorption}, answer={2.00}, answer_src=key, difficulty=5, diff_parts={K1 S4 I1 C0 F0}, diff_note={concepts: physisorption and chemisorption facts | steps: adsorbate is the substance adsorbed; easily liquefiable gases adsorb more; physisorption is reversible; Freundlich isotherm not linear at all pressures; charcoal in gas masks; mineral water meaning; heat of adsorption per mole; hydrogen on charcoal is physical; froth floatation uses adsorption}]
How many of the following statements is/are correct ?
\begin{qoptions}
\opt{1.}{Adsorbate is that substance where adsorption takes place.}
\opt{2.}{Lower the critical temperature of the gas, the more readily it will be adsorbed.}
\opt{3.}{Physical adsorption is an irreversible process.}
\opt{4.}{The amount of the gas adsorbed is directly proportional to the pressure, in all pressure region.}
\opt{5.}{By the gas masks, the poisonous gases are removed from air by charcoal, using process of adsorption.}
\opt{6.}{The water freed from all ions (cations and anions) is referred as mineral water.}
\opt{7.}{Heat of adsorption is defined as the energy liberated when $1$ gram of a gas is adsorbed on the solid surface.}
\opt{8.}{The adsorption of hydrogen on charcoal is chemical adsorption.}
\opt{9.}{Froth floatation process is an application of adsorption.}
\end{qoptions}
\end{bankq}
\begin{bankq}[id=JA04-P1-CHEM-S2-Q07, source={JEE Advanced (2021 Series) - Test 04 (04 Apr 2021)}, paper=1, qno=7, type=multi, subtopic={Colloids}, answer={A,C}, answer_src=key, difficulty=4, diff_parts={K1 S2 I1 C0 F1}, diff_note={concepts: properties of lyophilic and lyophobic colloids | steps: lyophilic sols are reversible while lyophobic are irreversible; lyophilic dispersed phase has high affinity for the medium; lyophilic particles are macromolecules in colloidal range; lyophobic particles are aggregates near the upper limit}]
Which of the following is/are correct statement(s).
\begin{qoptions}
\opt{A}{Lyophilic solution are reversible whereas lyophobic solution are irreversible}
\opt{B}{Dipersed phase shows no affinity with dispersion medium in case of lyophilic colloidal}
\opt{C}{Dipersed phase molecules of lyophilic colloidals are large enough to be in colloidal range}
\opt{D}{Dispersed phase molecules of lyophobic colloidals are large enough to be close to upper limit of colloidal range}
\end{qoptions}
\end{bankq}
\begin{bankq}[id=JA08-P2-CHEM-S1-Q06, source={JEE Advanced (2021 Series) - Test 08 (01 May 2021)}, paper=2, qno=6, type=single, subtopic={Colloids}, answer={C}, answer_src=solved, working={haemoglobin is a positively charged colloid, so statement C is incorrect.}, difficulty=3, diff_parts={K1 S2 I1 C0 F0}, diff_note={concepts: stability and charge of colloids | steps: Brownian motion keeps particles suspended; true solutions show no Tyndall effect; haemoglobin sol is positively charged while gold sol is negative; casein stabilises milk emulsion}]
Which is incorrect statement ?
\begin{qoptions}
\opt{A}{Brownian movement is responsible for the stability of colloidal solution}
\opt{B}{No Tyndall effect observed in True solution}
\opt{C}{Haemoglobin (blood) and gold sol are negative charge colloids}
\opt{D}{Casein act as emulsifier for milk}
\end{qoptions}
\end{bankq}
\begin{bankq}[id=JA11-P2-CHEM-S1-Q06, source={JEE Advanced (2021 Series) - Test 11 (20 May 2021)}, paper=2, qno=6, type=single, subtopic={Colloids}, answer={B}, answer_src=key, difficulty=2, diff_parts={K1 S2 I0 C0 F0}, diff_note={concepts: properties of colloids and emulsions | steps: milk is an emulsion of fat in water; lyophilic sols lower the surface tension of water slightly while lyophobic sols have no effect; conductivity is not negligible for lyophilic sols; option B}, also_in={JEE Advanced (2021 Series) - Test 15 (12 Jun 2021) P2 Q6; Nurture (General) - Test 01 (12 Jun 2021) P2 Q6}]
Which of the following statements is correct ?
\begin{qoptions}
\opt{A}{Milk is not an emulsion.}
\opt{B}{Lyophilic hydrosols have lower surface tension than H$_2$O.}
\opt{C}{Lyophobhic hydrosols have lower surface tension than water.}
\opt{D}{Conductivity of Lyophilic colloids is negligible}
\end{qoptions}
\end{bankq}
\begin{bankq}[id=JA12-P1-CHEM-S1-Q04, source={JEE Advanced (2021 Series) - Test 12 (25 May 2021)}, paper=1, qno=4, type=multi, subtopic={Colloids}, answer={A,B,C}, answer_src=solved, working={Smoke, ruby glass and pumice stone are colloidal systems; chlorophyll in solution is not.}, difficulty=3, diff_parts={K1 S2 I0 C0 F1}, diff_note={concepts: dispersed systems that are colloids | steps: smoke is a solid in gas; ruby glass is a solid sol of gold; pumice is a solid foam of gas in solid; chlorophyll is dissolved in molecular form}]
Which of the following is/are colloids ?
\begin{qoptions}
\opt{A}{Smoke}
\opt{B}{Ruby glass}
\opt{C}{Pumice stone}
\opt{D}{Chlorophyll}
\end{qoptions}
\end{bankq}
\begin{bankq}[id=JA13-P1-CHEM-S2-Q04, source={JEE Advanced (2021 Series) - Test 13 (31 May 2021)}, paper=1, qno=4, type=single, subtopic={Colloids}, answer={A}, answer_src=solved, working={Peptization produces a colloid and involves adsorption of common ions.}, difficulty=1, diff_parts={K1 S1 I0 C0 F0}, diff_note={concepts: peptization and common ion adsorption | steps: peptization converts a precipitate to a sol by adsorption of a common ion and is a preparation technique}]
Which of the following combination is correct?
\begin{qoptions}
\opt{A}{(I) (i) (S)}
\opt{B}{(II) (i) (R)}
\opt{C}{(III) (ii) (S)}
\opt{D}{(I) (ii) (S)}
\end{qoptions}
\end{bankq}
\begin{bankq}[id=JA13-P1-CHEM-S2-Q05, source={JEE Advanced (2021 Series) - Test 13 (31 May 2021)}, paper=1, qno=5, type=single, subtopic={Colloids}, answer={B}, answer_src=solved, working={Delta formation is coagulation (destabilisation).}, difficulty=1, diff_parts={K1 S1 I0 C0 F0}, diff_note={concepts: coagulation of colloids at river deltas | steps: delta formation occurs by coagulation of sol particles by electrolytes in sea water which is destabilisation}]
Which of the following combination is correct?
\begin{qoptions}
\opt{A}{(II) (i) (S)}
\opt{B}{(II) (ii) (R)}
\opt{C}{(III) (ii) (S)}
\opt{D}{(IV) (iii) (R)}
\end{qoptions}
\end{bankq}
\begin{bankq}[id=JA13-P1-CHEM-S2-Q06, source={JEE Advanced (2021 Series) - Test 13 (31 May 2021)}, paper=1, qno=6, type=single, subtopic={Colloids}, answer={D}, answer_src=solved, working={Peptization is not destabilisation.}, difficulty=3, diff_parts={K1 S2 I0 C1 F0}, diff_note={concepts: properties of colloids | steps: Tyndall effect is optical and the sky colour is scattering; delta formation is destabilisation by coagulation; protective power is measured by the gold number and is stabilisation; peptization is not destabilisation so one option is incorrect | C: four combinations to test}]
Which of the following combination is INCORRECT?
\begin{qoptions}
\opt{A}{(III) (iv) (P)}
\opt{B}{(II) (ii) (R)}
\opt{C}{(IV) (iii) (Q)}
\opt{D}{(I) (ii) (Q)}
\end{qoptions}
\end{bankq}
\begin{bankq}[id=JA14-P1-CHEM-S2-Q01, source={JEE Advanced (2021 Series) - Test 14 (06 Jun 2021)}, paper=1, qno=1, type=multi, subtopic={Adsorption}, answer={A,B,D}, answer_src=solved, working={Chemisorption is irreversible, specific, monolayer with heat of adsorption 80 to 240 kJ.}, difficulty=2, diff_parts={K1 S1 I0 C0 F1}, diff_note={concepts: characteristics of chemisorption | steps: chemisorption is irreversible, specific, monolayer with high enthalpy}]
Which of the following is/are characteristic of chemisorptions
\begin{qoptions}
\opt{A}{It is irreversible}
\opt{B}{It is specific}
\opt{C}{It is multilayer phenomenon}
\opt{D}{Magnitude of heat of adsorption is 80-240 kJ}
\end{qoptions}
\end{bankq}
\begin{bankq}[id=JA14-P2-CHEM-S3-Q04, source={JEE Advanced (2021 Series) - Test 14 (06 Jun 2021)}, paper=2, qno=4, type=single, subtopic={Adsorption}, answer={5}, answer_src=key, working={Printed key.}, difficulty=4, diff_parts={K1 S2 I1 C1 F0}, diff_note={concepts: surface area from monolayer adsorption volume | steps: moles of nitrogen from 112 cm3 at STP; molecules from Avogadro; area per molecule times number; convert units to m2 | I: the monolayer coverage links volume to area | C: unit conversion across cm3 and m2}]
The volume of nitrogen gas (at STP) required to cover a sample of silica gel with a monomolecular layer is 112 cm$^3$g$^{-1}$ of gel. Calculate the surface area per gram (in m$^2$) of the gel if each nitrogen molecule occupies $16.0\times10^{-22}$m$^2$. Give your answer to the nearest single digit integer $\{N_A=6\times10^{23}\}$.
\end{bankq}

\section*{TOAS papers}
\begin{bankq}[id=TOAS17-P2-CHEM-S2-Q10, source={TOAS - Test 17 (18 Sep 2022)}, paper=2, qno=10, type=single, subtopic={Catalysis}, answer={321}, answer_src=solved, working={V2O5 is the contact process catalyst (3), CuCl2 the Deacon process (2), Pt/Rh the Ostwald process (1)}, difficulty=1, diff_parts={K1 S1 I0 C0 F0}, diff_note={concepts: industrial catalysts | steps: vanadium pentoxide in the contact process; Deacon and Ostwald catalysts}]
\begin{center}\footnotesize\setlength{\tabcolsep}{2pt}
\begin{tabular}{|L{0.6cm}L{1.5cm}|L{0.6cm}L{2.2cm}|}\hline
\multicolumn{2}{|c|}{List-I Catalyst}&\multicolumn{2}{c|}{List-II Process}\\\hline
(P)&V$_2$O$_5$&(1)&Ostwald process\\\hline
(Q)&CuCl$_2$&(2)&Deacon process\\\hline
(R)&Pt/Rh&(3)&Contact process\\\hline
\end{tabular}\end{center}
Match list-I and list-II and give your answer as abc where a,b and c are the match of 1, 2, 3 from list-II.

\textbf{(for example Matching of P is = 2, Q = 3 \& R = 1 then your answer is 231)}
\end{bankq}
\begin{bankq}[id=TOAS18-P2-CHEM-S2-Q08, source={TOAS - Test 18 (30 Oct 2022)}, paper=2, qno=8, type=multi, subtopic={Colloids}, answer={B,C}, answer_src=key, difficulty=3, diff_parts={K1 S2 I0 C0 F1}, diff_note={concepts: properties of colloids | steps: peptisation; Tyndall effect; micelles}]
The correct statement(s) related to colloids is(are)
\begin{qoptions}
\opt{A}{The process of precipitating colloidal sol by adding an electrolyte is called peptization.}
\opt{B}{All colloids may not show Tyndall effect significantly.}
\opt{C}{Surfactants form micelle above critical micelle concentration (CMC), if temperature is above Kraft temperature.}
\opt{D}{Micelles are macromolecular colloids.}
\end{qoptions}
\end{bankq}
\end{document}
